\documentclass[11pt,a4paper]{ctexart}
% ============================================================
%  用极限定义证明 · 专题讲义
%  编译：xelatex 两遍（需在同一目录）
%  依据：上海交大高数 A 类期中真题 2008--2025 共 17 套（第 11 题逐卷核对）
%        官方答案原文经 4x 渲染逐字读取；本文全部 delta / X 常数已数值复核
%        （Homework/_review/tools/verify_handout_examples.py）
% ============================================================
\usepackage[a4paper,top=2.3cm,bottom=2.3cm,left=2.2cm,right=2.2cm]{geometry}
\usepackage{amsmath,amssymb,amsthm,mathtools,bm}
\usepackage{booktabs,array,longtable,tabularx,multicol}
\usepackage{enumitem}
\usepackage[most]{tcolorbox}
\usepackage{fancyhdr}
\usepackage{hyperref}
\hypersetup{hidelinks,colorlinks=true,linkcolor=blue!50!black,urlcolor=blue!50!black}
\usepackage{titlesec}
\usepackage{pifont}

% ---------- 版式 ----------
\linespread{1.28}
\setlength{\parskip}{0.25em}
\setlist[enumerate]{itemsep=3pt,topsep=4pt,parsep=0pt,leftmargin=1.9em}
\setlist[itemize]{itemsep=3pt,topsep=4pt,parsep=0pt,leftmargin=1.6em}
\titleformat{\section}{\Large\bfseries\color{blue!35!black}}{\thesection}{0.6em}{}
\titleformat{\subsection}{\large\bfseries\color{blue!25!black}}{\thesubsection}{0.5em}{}
\titlespacing*{\section}{0pt}{1.7ex plus .3ex}{1.1ex}
\titlespacing*{\subsection}{0pt}{1.3ex plus .2ex}{0.8ex}
\setcounter{tocdepth}{2}

\pagestyle{fancy}
\fancyhf{}
\fancyhead[L]{\small 用极限定义证明 · 专题讲义}
\fancyhead[R]{\small 陶理}
\fancyfoot[C]{\small\thepage}
\renewcommand{\headrulewidth}{0.4pt}

% ---------- 数学简写 ----------
\newcommand{\R}{\mathbb{R}}
\newcommand{\Z}{\mathbb{Z}}
\newcommand{\eps}{\varepsilon}
\newcommand{\abs}[1]{\left\lvert #1\right\rvert}
\newcommand{\ee}{\mathrm{e}}
\renewcommand{\proofname}{证明}

% ---------- 盒子 ----------
\newtcolorbox{keypoint}[1][]{enhanced,breakable,colback=blue!3,colframe=blue!40!black,
  fonttitle=\bfseries,title={#1},boxrule=0.5pt,arc=1.5mm,
  left=3mm,right=3mm,top=2mm,bottom=2mm,before skip=4pt,after skip=4pt}
\newtcolorbox{warnbox}[1][]{enhanced,breakable,colback=red!3,colframe=red!50!black,
  fonttitle=\bfseries,title={#1},boxrule=0.5pt,arc=1.5mm,
  left=3mm,right=3mm,top=2mm,bottom=2mm,before skip=4pt,after skip=4pt}
\newtcolorbox{methodbox}[1][]{enhanced,breakable,colback=green!3!white,colframe=green!35!black,
  fonttitle=\bfseries,title={#1},boxrule=0.5pt,arc=1.5mm,
  left=3mm,right=3mm,top=2mm,bottom=2mm,before skip=4pt,after skip=4pt}
\newtcolorbox{notebox}[1][]{enhanced,breakable,colback=gray!5,colframe=gray!55,
  fonttitle=\bfseries,title={#1},boxrule=0.4pt,arc=1mm,
  left=3mm,right=3mm,top=1.5mm,bottom=1.5mm,before skip=3pt,after skip=3pt}

% ---------- 例题 ----------
\newcounter{exa}[section]
\newenvironment{exa}[1][]{\refstepcounter{exa}\par\smallskip\noindent
  \textbf{例 \thesection.\theexa}\if\relax\detokenize{#1}\relax\else\textbf{（#1）}\fi\ \ignorespaces}{\par\smallskip}
\newenvironment{solution}{\par\smallskip\noindent\textbf{解：}\ignorespaces}{\par\smallskip}
\newenvironment{analysis}{\par\smallskip\noindent\textbf{分析：}\ignorespaces}{\par\smallskip}

\newcommand{\kw}[1]{\textbf{\textcolor{blue!45!black}{#1}}}
\newcommand{\yes}{\textcolor{green!50!black}{\ding{51}}}
\newcommand{\no}{\textcolor{red!70!black}{\ding{55}}}

\begin{document}

% ============================================================
%  封面
% ============================================================
\begin{titlepage}
\centering
\vspace*{2.0cm}
{\Huge\bfseries 用极限定义证明\par}
\vspace{0.45cm}
{\Huge\bfseries 专题讲义\par}
\vspace{0.9cm}
{\large 三种形式 \quad$\cdot$\quad 一套骨架 \quad$\cdot$\quad 十道真题精讲\par}
\vspace{1.4cm}
{\large 依据 \textbf{上海交通大学 2008--2025 级第一学期高数 A 类期中真题 17 套}\par}
\vspace{0.3cm}
{\large 第 11 题逐卷核对（15 套卷面 + 2 套仅答案）\par}
\vspace{0.3cm}
{\large 官方答案经 4 倍渲染逐字读取，$\delta$ / $X$ 常数全部数值复核\par}
\vspace{2.0cm}
{\large 陶理\ \ 电子信息与电气工程学院\par}
\vspace{0.3cm}
{\large 2026 年 11 月\par}
\vfill
{\small\begin{minipage}{0.86\textwidth}
\textbf{怎么用这份讲义}\par\vspace{2pt}
第 1--2 节：搞清楚定义到底在说什么（含量词顺序与否定式）；\quad
第 3 节：五步骨架与“先关区间”的理由；\quad
第 4 节：十道真题精讲，\textbf{先遮住解答自己做}; \quad
第 5 节：十条高频错误对照；\quad
第 6 节：六道练习（附答案）；\quad
附录 A：17 套真题全表（含每题的关键 $\delta$ / $X$）；\quad
附录 B：考前 20 分钟速查。\par\vspace{3pt}
\textbf{配套}：知识库速查版 \texttt{Note/Advanced\_Mathematics/Notes/}（极限定义证明 · 题型专题）；数值复核脚本 \texttt{Homework/\_review/tools/verify\_handout\_examples.py}。
\end{minipage}}
\end{titlepage}

\tableofcontents
\newpage

% ============================================================
\section{为什么需要 $\eps$--$\delta$：从“无限接近”到“可以检验”}
% ============================================================

“$x$ 无限接近 $x_0$ 时 $f(x)$ 无限接近 $A$”这句话在直觉上够用，但它有两个致命问题：\textbf{“无限接近”本身没有定义}，而且\textbf{它不能被检验}。数学家把这句话翻译成了下面这个可以逐条核对的句式：

\begin{keypoint}[极限的定义（$\eps$--$\delta$ 型）]
设 $f$ 在 $x_0$ 的某去心邻域内有定义。若存在常数 $A$，使得
\[
\forall \eps>0,\ \exists \delta>0,\ \forall x:\quad 0<\abs{x-x_0}<\delta \ \Longrightarrow\ \abs{f(x)-A}<\eps,
\]
则称 $\lim_{x\to x_0}f(x)=A$。
\end{keypoint}

把它读成一场\textbf{挑战--应答博弈}，就不会再觉得抽象：

\begin{center}
\small
\begin{tabular}{@{}clll@{}}
\toprule
回合 & 谁出牌 & 出什么 & 约束 \\
\midrule
1 & 考官 & 精度 $\eps>0$ & 可以任意小，但一旦给出就固定 \\
2 & 你 & 阈值 $\delta>0$ & \textbf{只能依赖 $\eps$}（和题中的常数），不能依赖 $x$ \\
3 & 考官 & 一个点 $x$ & 只要求 $0<\abs{x-x_0}<\delta$（可以任意刁钻） \\
4 & 判定 & 是否 $\abs{f(x)-A}<\eps$ & 必须成立，一个反例就输 \\
\bottomrule
\end{tabular}
\end{center}

\begin{methodbox}[一句话记住]
\textbf{$\eps$ 是“允许多大的误差”，$\delta$ 是“离多近才能保证这个误差”。}
证明的全部工作，就是\textbf{从 $\eps$ 反解出一个只依赖 $\eps$ 的 $\delta$}。
\end{methodbox}

\subsection{三个常见的“以为懂了”}

\begin{enumerate}
\item \textbf{“$\delta$ 越小越保险”}——错。$\delta$ 是反解出来的：写小了不违反定义，但拿不到分步；写大了结论不成立。阅卷看的是“你的 $\delta$ 是否恰好让不等式成立”。
\item \textbf{“$\delta$ 和 $\eps$ 无关”}——错。$\delta$ 必须随 $\eps$ 变小而变小（通常 $\delta\sim\eps/C$）。这正是“连续性/收敛性”能被量化的根源。
\item \textbf{“随便取一个很小的 $\delta$ 就行”}——错。$\delta$ 一旦固定，考官可以挑区间里最坏的那个 $x$。所以你必须\textbf{对区间内所有 $x$ 一致地}控制住 $\abs{f(x)-A}$。
\end{enumerate}

% ============================================================
\section{三种定义：同一件事的三副面孔}
% ============================================================

\subsection{模板 A：$\eps$--$\delta$ 型（$x\to x_0$，有限点）}

\begin{keypoint}[定义 A]
\[
\lim_{x\to x_0}f(x)=A
\iff
\forall \eps>0,\ \exists \delta>0,\ \forall x:\quad 0<\abs{x-x_0}<\delta \Rightarrow \abs{f(x)-A}<\eps .
\]
\end{keypoint}

两个细节：$0<\abs{x-x_0}$ 是\textbf{去心}（$x=x_0$ 处函数值不影响极限）；$\abs{x-x_0}<\delta$ 同时管左右两侧，所以\textbf{双侧极限}要求左右同时成立。

\textbf{单侧版本}：把 $0<\abs{x-x_0}<\delta$ 换成 $x_0<x<x_0+\delta$（右极限）或 $x_0-\delta<x<x_0$（左极限）。

\subsection{模板 B：$\eps$--$X$ 型（$x\to+\infty$ 或 $x\to-\infty$）}

\begin{keypoint}[定义 B]
\[
\lim_{x\to+\infty}f(x)=A
\iff
\forall \eps>0,\ \exists X>0,\ \forall x:\quad x>X \Rightarrow \abs{f(x)-A}<\eps .
\]
$x\to-\infty$ 时把 $x>X$ 换成 $x<-X$；$x\to\infty$ 时换成 $\abs{x}>X$。
\end{keypoint}

与模板 A 的差别只有两处：把“离 $x_0$ 多近”换成“$x$ 有多大”，把“去心邻域”换成“一条半直线”。\textbf{阈值符号从 $\delta$ 变成 $X$}，收尾从 $\min$ 变成 $\max$。

\subsection{模板 C：无穷大 $M$ 型（结论是 $\infty$）}

\begin{keypoint}[定义 C]
\[
f(x)\to+\infty\ (x\to+\infty)
\iff
\forall M>0,\ \exists X>0,\ \forall x:\quad x>X \Rightarrow f(x)>M .
\]
\[
\abs{f(x)}\to\infty\ (x\to x_0)
\iff
\forall M>0,\ \exists \delta>0,\ \forall x:\quad 0<\abs{x-x_0}<\delta \Rightarrow \abs{f(x)}>M .
\]
\end{keypoint}

\begin{warnbox}[方向和 A/B 相反]
A/B 要的是“误差\textbf{小}”：把 $\abs{f-A}$ 往\textbf{大了估}（上界），再令上界 $<\eps$。\\
C 要的是“函数值\textbf{大}”：把 $\abs{f}$ 往\textbf{小了估}（下界），再令下界 $>M$。\\
\textbf{方向搞反，是这一题型最致命的错}（连官方答案都栽过，见例 4.9、4.10）。
\end{warnbox}

\subsection{三种形式对照表}

\begin{center}
\small
\begin{tabular}{@{}llll@{}}
\toprule
 & 模板 A（$\eps$--$\delta$） & 模板 B（$\eps$--$X$） & 模板 C（无穷大 $M$） \\
\midrule
任意给定量 & $\eps>0$ & $\eps>0$ & $M>0$ \\
存在量 & $\delta>0$ & $X>0$ & $\delta>0$ 或 $X>0$ \\
适用范围 & $0<\abs{x-x_0}<\delta$ & $x>X$（或 $\abs{x}>X$） & $0<\abs{x-x_0}<\delta$ 或 $x>X$ \\
要证的不等式 & $\abs{f-A}<\eps$ & $\abs{f-A}<\eps$ & $\abs{f}>M$（或 $f>M$） \\
放缩目标 & $C\abs{x-x_0}$ & $C/\abs{x}$ & $k/\abs{x-x_0}$ 或 $k\,g(x)$ \\
收尾 & $\delta=\min\{\delta_1,\eps/C\}$ & $X=\max\{X_1,C/\eps\}$ & $\delta=\min\{\delta_1,k/M\}$ \\
\bottomrule
\end{tabular}
\end{center}

\subsection{量词顺序为什么不能换}

定义里的 $\forall\eps\,\exists\delta\,\forall x$ 是一个整体。换掉任何两个量词，含义都会变：

\begin{itemize}
\item $\exists\delta\,\forall\eps$：存在一个 $\delta$ 对所有 $\eps$ 都管用。取 $\eps\to0^+$ 立刻推出 $f$ 在该邻域内\textbf{恒等于} $A$——这比“极限存在”强得多，绝大多数函数不满足。
\item $\forall x\,\exists\delta$：$\delta$ 可以依赖 $x$。那每个人都能取“离自己足够近”的邻域，命题几乎恒真，等于什么都没说。
\item $\forall\eps\,\forall\delta$：任何 $\delta$ 都行——显然过强。
\end{itemize}

\begin{notebox}[考场提醒]
写定义时\textbf{按顺序默写} “任给 $\eps>0$，存在 $\delta>0$，当 $0<\abs{x-x_0}<\delta$ 时，有 $\abs{f(x)-A}<\eps$”，
不要跳步，也不要把“$\exists\delta$”写成“取 $\delta=\cdots$”（那是求解阶段的事）。
\end{notebox}

\subsection{否定式：证明极限不存在}

把量词依次取反、结论取反，就得到“$A$ 不是极限”的严格含义：
\[
\exists \eps_0>0,\ \forall \delta>0,\ \exists x:\quad 0<\abs{x-x_0}<\delta \ \text{且}\ \abs{f(x)-A}\ge \eps_0 .
\]
注意 $\eps_0$ 是\textbf{先}给定的（可以自己挑一个方便的），而 $x$ 可以依赖 $\delta$。用法：对任意候选 $A$，找一个 $\eps_0$，使得无论 $\delta$ 多小，区间里总有“坏点”。

\begin{exa}[否定式的典型用法]
证明 $\lim_{x\to0}\sin\dfrac1x$ 不存在。
\end{exa}
\begin{solution}
反设极限为 $A$。取 $\eps_0=\dfrac12$。对任意 $\delta>0$，取正整数 $n$ 充分大，使
\[
x_n=\frac{1}{\frac{\pi}{2}+2n\pi}\in(0,\delta),\qquad
y_n=\frac{1}{\frac{3\pi}{2}+2n\pi}\in(0,\delta),
\]
则 $\sin\frac1{x_n}=1$，$\sin\frac1{y_n}=-1$。若 $\abs{f(x_n)-A}<\frac12$ 且 $\abs{f(y_n)-A}<\frac12$ 同时成立，则
\[
2=\abs{f(x_n)-f(y_n)}\le\abs{f(x_n)-A}+\abs{f(y_n)-A}<1,
\]
矛盾。所以至少有一个“坏点”，即不存在这样的 $A$。$\blacksquare$
\end{solution}

% ============================================================
\section{五步骨架：所有真题共用一套流程}
% ============================================================

\begin{methodbox}[考场只背这五步]
\begin{enumerate}
\item \kw{写定义}：把题面翻译成模板 A / B / C 的 $\forall$--$\exists$ 句式（先写出来，后面每一行都服务于它）。
\item \kw{化差}：把 $\abs{f(x)-A}$ 变成能控制的形式——有理化、通分、拆项、乘共轭、三角不等式。
\item \kw{先关区间}：用一个 $\delta_1$（或 $X_1$）把 $x$ 限制在小范围，使所有分母有\textbf{正下界}、根号内有\textbf{正下界}、放缩常数\textbf{定死}。
\item \kw{放缩}：得到 $C\abs{x-x_0}<\eps$（或 $C/\abs{x}<\eps$、$k/\abs{x-x_0}>M$）。
\item \kw{反解}：$\delta=\min\{\delta_1,\eps/C\}$，或 $X=\max\{X_1,C/\eps\}$，或 $\delta=\min\{\delta_1,k/M\}$。
\end{enumerate}
\end{methodbox}

三种形式\textbf{只差第 1 步和第 5 步}，第 2--4 步完全通用——这就是“背一套骨架、应付三种形式”的底气。

\subsection{第 3 步为什么不能省}

看一个具体例子。要证 $\lim_{x\to2}\dfrac{1}{x-1}=1$，化差得
\[
\abs{\frac{1}{x-1}-1}=\frac{\abs{x-2}}{\abs{x-1}}.
\]
如果\textbf{不限制} $x$ 的范围，那么当 $x$ 靠近 $1$ 时 $\abs{x-1}$ 可以任意小，$\dfrac{1}{\abs{x-1}}$ 无上界——\textbf{找不到一个对所有 $x$ 都成立的常数 $C$}。而一旦取 $\delta_1=\dfrac12$，就有 $x\in\left(\dfrac32,\dfrac52\right)$，从而 $\abs{x-1}>\dfrac12$，常数 $C=2$ 立刻定死。

\begin{keypoint}[一句话]
\textbf{“先关区间”不是为了严谨好看，而是为了让放缩常数存在。}
没有第 3 步，第 4 步的常数就是无源之水。
\end{keypoint}

\subsection{$\delta_1$ 的常用取法}

\begin{center}
\small
\begin{tabular}{@{}lll@{}}
\toprule
目的 & 取法 & 说明 \\
\midrule
让分母远离 $0$ & $\delta_1=\dfrac12$ 或 $\dfrac{\abs{x_0}}{2}$ & 保证 $\abs{x}\ge\abs{x_0}/2>0$ \\
让 $\abs{x}\ge1$ & $\delta_1=\abs{x_0}-1$（需 $>0$） & 便于把 $\abs{x}$ 换成常数 \\
让根号内有下界 & $\delta_1$ 取到“根号内 $\ge c>0$” & 如 $\sqrt{3-x}$ 需 $x\le 3-\eta$ \\
$x\to+\infty$ 时让 $x\ge1$ & $X_1=1$ 或 $2$ & 之后可用 $x^2/2$ 型估计 \\
含参数 $a$ & $X_1=\max\{1,\sqrt{2(\abs{a}+1)}\}$ & 把 $\abs{a}$ 压成常数 \\
\bottomrule
\end{tabular}
\end{center}

\begin{warnbox}[收尾必须是 $\min$ / $\max$]
$\delta_1$ 管“区间”，$\eps/C$ 管“精度”，两个条件必须\textbf{同时}满足，所以 $\delta=\min\{\delta_1,\eps/C\}$。
只写 $\delta=\eps/C$ 而忘了 $\delta_1$，等于默认 $\eps/C\le\delta_1$——$\eps$ 一大就不成立。
类似地，$X=\max\{X_1,C/\eps\}$ 而不是 $\min$。
\end{warnbox}

\subsection{放缩的三条硬规则}

\begin{enumerate}
\item \textbf{先关区间，再谈放缩}——否则常数定不下来。
\item \textbf{求上界时}：分子取上界、分母取下界；\textbf{求下界时}：分子取下界、分母取上界。方向反了整题报废。
\item \textbf{每写一个不等号，问一句“方向对吗、等号能取到吗”}——官方答案也栽在这一步。
\end{enumerate}

% ============================================================
\section{十道真题精讲}
% ============================================================

本节按模板分组：4.1--4.3 是 $\eps$--$\delta$，4.4--4.7 是 $\eps$--$X$，4.8--4.10 是无穷大 $M$ 型。\\[2pt]
\textbf{建议用法}：先遮住解答，自己写 8 分钟；再对照“分析”检查卡在哪一步。

\subsection{模板 A：$\eps$--$\delta$}

\begin{exa}[2011--2012 卷第 11 题]
用极限定义证明 $\displaystyle\lim_{x\to2}\frac{1}{x-1}=1$。
\end{exa}
\begin{solution}
\textbf{（1）写定义}：任给 $\eps>0$，需找 $\delta>0$，使 $0<\abs{x-2}<\delta$ 时 $\abs{\dfrac{1}{x-1}-1}<\eps$。

\textbf{（2）化差}：
\[
\abs{\frac{1}{x-1}-1}=\abs{\frac{2-x}{x-1}}=\frac{\abs{x-2}}{\abs{x-1}}.
\]
\textbf{（3）关区间}：取 $\delta_1=\dfrac12$，则 $0<\abs{x-2}<\delta_1$ 蕴含 $\dfrac32<x<\dfrac52$，故 $\abs{x-1}>\dfrac12$，于是
\[
\abs{\frac{1}{x-1}-1}<\frac{\abs{x-2}}{1/2}=2\abs{x-2}.
\]
\textbf{（4）放缩}：已得 $C=2$。

\textbf{（5）反解}：取 $\delta=\min\left\{\dfrac12,\dfrac{\eps}{2}\right\}$。则当 $0<\abs{x-2}<\delta$ 时，
\[
\abs{\frac{1}{x-1}-1}<2\abs{x-2}<2\delta\le \eps.\qquad\blacksquare
\]
\end{solution}

\begin{notebox}[检查自己的答案]
三个量必须对得上：(1) 区间内 $\abs{x-1}$ 的\textbf{下界} $\frac12$；(2) 放缩系数 $C=2$；(3) $\delta$ 里 $\frac12$ 和 $\frac{\eps}{2}$ 各司其职（一个管区间、一个管精度）。
\end{notebox}

\begin{exa}[2017--2018 卷第 11 题]
用极限定义证明 $\displaystyle\lim_{x\to1}x\sqrt{x+3}=2$。
\end{exa}
\begin{solution}
\textbf{（1）写定义}：任给 $\eps>0$，需找 $\delta>0$，使 $0<\abs{x-1}<\delta$ 时 $\abs{x\sqrt{x+3}-2}<\eps$。

\textbf{（2）化差（拆项）}：关键是把 $\sqrt{x+3}\to2$ 与 $x\to1$ 分开：
\[
x\sqrt{x+3}-2=x\bigl(\sqrt{x+3}-2\bigr)+2(x-1).
\]
而
\[
x\bigl(\sqrt{x+3}-2\bigr)=x\cdot\frac{(x+3)-4}{\sqrt{x+3}+2}=\frac{x(x-1)}{\sqrt{x+3}+2}.
\]
故
\[
\abs{x\sqrt{x+3}-2}\le \frac{\abs{x}}{\sqrt{x+3}+2}\abs{x-1}+2\abs{x-1}.
\]
\textbf{（3）关区间}：取 $\delta_1=1$，则 $0<x<2$，故 $\abs{x}<2$，且 $\sqrt{x+3}+2>2$，于是
\[
\frac{\abs{x}}{\sqrt{x+3}+2}<\frac{2}{2}=1.
\]
\textbf{（4）放缩}：$\abs{x\sqrt{x+3}-2}\le (1+2)\abs{x-1}=3\abs{x-1}$。

\textbf{（5）反解}：取 $\delta=\min\left\{1,\dfrac{\eps}{3}\right\}$，则 $0<\abs{x-1}<\delta$ 时上式 $<3\delta\le\eps$。$\blacksquare$
\end{solution}

\begin{notebox}[点评]
这道题的教学价值在于\textbf{拆项}。遇到“乘积型”$u(x)v(x)\to u_0v_0$，标准动作是
\[
uv-u_0v_0=u\,(v-v_0)+v_0\,(u-u_0).
\]
本题 $u=x$、$v=\sqrt{x+3}$、$u_0=1$、$v_0=2$，所以拆成 $x\bigl(\sqrt{x+3}-2\bigr)+2(x-1)$：前半段归到 $\sqrt{x+3}\to2$，后半段归到 $x\to1$，各自都能写成 $C\abs{x-1}$。
\end{notebox}

\begin{exa}[2021--2022 卷第 11 题]
利用极限定义证明 $\displaystyle\lim_{x\to1}\frac{\sqrt{3-x}-\sqrt{x+1}}{x-1}=-\frac{1}{\sqrt2}$。
\end{exa}
\begin{solution}
\textbf{（1）写定义}：任给 $\eps>0$，需找 $\delta>0$，使 $0<\abs{x-1}<\delta$ 时
$\abs{\dfrac{\sqrt{3-x}-\sqrt{x+1}}{x-1}+\dfrac{1}{\sqrt2}}<\eps$。

\textbf{（2）化差 · 第一层有理化}：分子分母同乘共轭，
\[
\frac{\sqrt{3-x}-\sqrt{x+1}}{x-1}
=\frac{(3-x)-(x+1)}{(x-1)\bigl(\sqrt{3-x}+\sqrt{x+1}\bigr)}
=-\frac{2}{S},\qquad S:=\sqrt{3-x}+\sqrt{x+1}.
\]
目标变为
\[
\abs{-\frac2S+\frac1{\sqrt2}}=\frac{\abs{S-2\sqrt2}}{\sqrt2\,S}.
\]
\textbf{（3）化差 · 第二层有理化}：
\[
\abs{S-2\sqrt2}
\le\abs{\sqrt{3-x}-\sqrt2}+\abs{\sqrt{x+1}-\sqrt2}
\le\frac{\abs{x-1}}{\sqrt{3-x}+\sqrt2}+\frac{\abs{x-1}}{\sqrt{x+1}+\sqrt2}.
\]
\textbf{（4）关区间}：取 $\delta_1=\dfrac12$，则 $\dfrac52<3-x<\dfrac72$、$\dfrac32<x+1<\dfrac52$，故
\[
\sqrt{3-x}+\sqrt2>\frac52,\qquad \sqrt{x+1}+\sqrt2>\frac52,\qquad S>\sqrt{\frac52}+\sqrt{\frac32}>2.
\]
于是
\[
\abs{S-2\sqrt2}<\frac{2\abs{x-1}}{5/2}=\frac45\abs{x-1},\qquad \sqrt2\,S>2\sqrt2>2.
\]
\textbf{（5）放缩}：
\[
\abs{-\frac2S+\frac1{\sqrt2}}<\frac{(4/5)\abs{x-1}}{2}=\frac25\abs{x-1}.
\]
\textbf{（6）反解}：取 $\delta=\min\left\{\dfrac12,\dfrac{5\eps}{2}\right\}$，则 $0<\abs{x-1}<\delta$ 时
\[
\abs{\frac{\sqrt{3-x}-\sqrt{x+1}}{x-1}+\frac1{\sqrt2}}<\frac25\abs{x-1}<\frac25\delta\le\eps.\qquad\blacksquare
\]
\end{solution}

\begin{methodbox}[两层有理化是“根式差”的固定动作]
第一层把 $\dfrac{\text{根式差}}{\text{差}}$ 化成 $\dfrac{\text{常数}}{\text{根式和}}$；第二层把 $\sqrt a-\sqrt b$ 换成 $\dfrac{a-b}{\sqrt a+\sqrt b}$。两次之后每一项都变成 $C\abs{x-x_0}$，可以统一放缩。
\end{methodbox}

\subsection{模板 B：$\eps$--$X$}

\begin{exa}[2010--2011 卷第 11 题]
用极限定义证明 $\displaystyle\lim_{x\to+\infty}\frac{x^2+x}{x^2-2}=1$。
\end{exa}
\begin{solution}
\textbf{（1）写定义}：任给 $\eps>0$，需找 $X>0$，使 $x>X$ 时 $\abs{\dfrac{x^2+x}{x^2-2}-1}<\eps$。

\textbf{（2）化差}：
\[
\abs{\frac{x^2+x}{x^2-2}-1}=\abs{\frac{x+2}{x^2-2}}=\frac{x+2}{x^2-2}\quad(x>2).
\]
\textbf{（3）关区间}：取 $X_1=2$，则 $x>2$ 时 $x^2>4$，故 $x^2-2>\dfrac{x^2}{2}$；同时 $x+2<2x$。

\textbf{（4）放缩}：
\[
\frac{x+2}{x^2-2}<\frac{2x}{x^2/2}=\frac4x .
\]
\textbf{（5）反解}：取 $X=\max\left\{2,\dfrac4\eps\right\}$，则 $x>X$ 时
\[
\abs{\frac{x^2+x}{x^2-2}-1}<\frac4x<\frac4X\le\eps.\qquad\blacksquare
\]
\end{solution}

\begin{exa}[2019--2020 卷第 11 题]
用极限定义证明 $\displaystyle\lim_{x\to\infty}\frac{x+1}{x^2+2}=0$。
\end{exa}
\begin{solution}
\textbf{（1）写定义}：任给 $\eps>0$，需找 $X>0$，使 $\abs{x}>X$ 时 $\abs{\dfrac{x+1}{x^2+2}}<\eps$。

\textbf{（2）化差}：目标 $A=0$，直接估计 $\abs{f}$：
\[
\abs{\frac{x+1}{x^2+2}}\le\frac{\abs{x}+1}{x^2+2}.
\]
\textbf{（3）关区间}：取 $X_1=1$，则 $\abs{x}>1$ 时 $\abs{x}+1\le2\abs{x}$，且 $x^2+2>x^2$。

\textbf{（4）放缩}：
\[
\abs{f(x)}\le\frac{2\abs{x}}{x^2}=\frac{2}{\abs{x}}.
\]
\textbf{（5）反解}：取 $X=\max\left\{1,\dfrac2\eps\right\}$，则 $\abs{x}>X$ 时 $\abs{f(x)}<\dfrac2{\abs{x}}<\dfrac2X\le\eps$。$\blacksquare$
\end{solution}

\begin{notebox}[点评]
注意“$x\to\infty$”是\textbf{双侧}的：$\abs{x}>X$。凡是出现 $\abs{x}$ 的估计（分子用 $\abs{x}$、分母用 $x^2$）天然两边通用；若题目是 $x\to-\infty$，绝对值估计同样适用，这正是模板 B 的优点。
\end{notebox}

\begin{exa}[2013--2014 卷第 11 题]
用极限定义证明 $\displaystyle\lim_{x\to+\infty}\frac{\sqrt{2x+\sqrt x}-\sqrt x}{1+x}=0$。
\end{exa}
\begin{solution}
\textbf{（1）写定义}：任给 $\eps>0$，需找 $X>0$，使 $x>X$ 时 $\abs{f(x)}<\eps$。

\textbf{（2）化差（有理化）}：
\[
\sqrt{2x+\sqrt x}-\sqrt x
=\frac{(2x+\sqrt x)-x}{\sqrt{2x+\sqrt x}+\sqrt x}
=\frac{x+\sqrt x}{\sqrt{2x+\sqrt x}+\sqrt x}.
\]
故
\[
f(x)=\frac{x+\sqrt x}{(1+x)\bigl(\sqrt{2x+\sqrt x}+\sqrt x\bigr)}.
\]
\textbf{（3）关区间}：取 $X_1=1$，则 $x>1$ 时 $x+\sqrt x\le 2x$，$1+x\ge x$，且 $\sqrt{2x+\sqrt x}+\sqrt x\ge\sqrt x$。

\textbf{（4）放缩}：
\[
\abs{f(x)}\le\frac{2x}{x\cdot\sqrt x}=\frac{2}{\sqrt x}.
\]
\textbf{（5）反解}：取 $X=\max\left\{1,\dfrac{4}{\eps^2}\right\}$（因为 $\dfrac{2}{\sqrt X}\le\eps\iff X\ge\dfrac{4}{\eps^2}$），则 $x>X$ 时 $\abs{f(x)}<\eps$。$\blacksquare$
\end{solution}

\begin{exa}[2018--2019 卷第 11 题（含参）]
设常数 $a\in\R$，用极限定义证明 $\displaystyle\lim_{x\to+\infty}\frac{x^2+\ee^{-x}}{x^3-ax+1}=0$。
\end{exa}
\begin{solution}
\textbf{（1）写定义}：任给 $\eps>0$，需找 $X>0$，使 $x>X$ 时 $\abs{f(x)}<\eps$。

\textbf{（2）化差}：$A=0$，直接估。分子 $x^2+\ee^{-x}\le x^2+1$（$x>0$ 时 $\ee^{-x}\le1$）；分母取\textbf{下界}：
\[
\abs{x^3-ax+1}\ge x^3-\abs{a}x-1 .
\]
\textbf{（3）关区间}：取
\[
X_1=\max\left\{1,\sqrt{2(\abs{a}+1)}\right\}.
\]
当 $x\ge X_1$ 时 $x\ge1$，故 $\dfrac1x\le1$，于是
\[
\frac{x^2}{2}\ge \abs{a}+1\ge \abs{a}+\frac1x
\ \Longrightarrow\ \frac{x^3}{2}\ge \abs{a}x+1
\ \Longrightarrow\ x^3-\abs{a}x-1\ge\frac{x^3}{2}.
\]
\textbf{（4）放缩}：$x\ge1$ 时 $x^2+1\le2x^2$，故
\[
\abs{f(x)}\le\frac{2x^2}{x^3/2}=\frac4x .
\]
\textbf{（5）反解}：取 $X=\max\left\{X_1,\dfrac4\eps\right\}$，则 $x>X$ 时 $\abs{f(x)}<\dfrac4X\le\eps$。$\blacksquare$
\end{solution}

\begin{methodbox}[含参题的通用套路]
参数一律先“控住”：用 $\abs{x^3-ax+1}\ge x^3-\abs{a}x-1$ 把 $a$ 的影响压成常数 $\abs{a}$，再用 $X_1$ 让分母的高阶项占优。\textbf{不要试图解关于 $a$ 的不等式}。
\end{methodbox}

\subsection{模板 C：无穷大 $M$ 型}

\begin{exa}[2009--2010 卷第 11 题]
用无穷大定义证明 $\displaystyle\lim_{x\to+\infty}(x^3-4x)=+\infty$。
\end{exa}
\begin{solution}
\textbf{（1）写定义}：任给 $M>0$，需找 $X>0$，使 $x>X$ 时 $x^3-4x>M$。

\textbf{（2）化差（估下界）}：要证“大”，所以往\textbf{小了估}。当 $x\ge3$ 时 $x^2\ge9>8$，故 $4x<\dfrac{x^3}{2}$，于是
\[
x^3-4x>\frac{x^3}{2}.
\]
\textbf{（3）关区间}：取 $X_1=3$。

\textbf{（4）放缩}：$\dfrac{x^3}{2}>M\iff x>(2M)^{1/3}$。

\textbf{（5）反解}：取 $X=\max\left\{3,(2M)^{1/3}\right\}$，则 $x>X$ 时
\[
x^3-4x>\frac{x^3}{2}>\frac{X^3}{2}\ge M.\qquad\blacksquare
\]
\end{solution}

\begin{warnbox}[官方答案有错 · 1]
\textbf{2015--2016 卷第 11 题}：$\displaystyle\lim_{x\to1}\frac{2-\arctan x}{1-x^2}=\infty$。

官方答案写
\[
\abs{\frac{2-\arctan x}{1-x^2}}>\frac{2-\pi/2}{\frac12\abs{1-x}}=\frac{k}{\abs{1-x}},\qquad k=4-\pi,
\]
从而 $\delta=\min\left\{\frac12,\frac{4-\pi}{M}\right\}$。\textbf{两处错误}：

\textbf{(1) 放缩方向错}：$x\to1$ 时 $\abs{1+x}\to2$，所以 $\abs{1+x}$ 需要的是\textbf{上界}。写成 $\frac12$ 等于要求 $\abs{1+x}<\frac12$，在 $x\to1$ 附近根本不成立。

\textbf{(2) 常数没跟着改}：正确的上界是 $\abs{1+x}<\frac52$（取 $\delta_1=\frac12$），常数应由 $4-\pi$ 改为 $\frac{4-\pi}{5}$。

\textbf{反例}：取 $M=10$，官方 $\delta=\frac{4-\pi}{10}\approx0.08584$；取 $x=1.07953$（满足 $0<\abs{x-1}<\delta$），实算 $\abs{f(x)}\approx7.11<10$——反例成立。该区间上 $\abs{f}$ 的最小值约 $6.55$（在 $x=1.08584$ 处），更不满足要求。
\end{warnbox}

\begin{exa}[2015--2016 卷第 11 题（订正版）]
用极限定义证明 $\displaystyle\lim_{x\to1}\frac{2-\arctan x}{1-x^2}=\infty$。
\end{exa}
\begin{solution}
\textbf{（1）写定义}：任给 $M>0$，需找 $\delta>0$，使 $0<\abs{x-1}<\delta$ 时 $\abs{\dfrac{2-\arctan x}{1-x^2}}>M$。

\textbf{（2）化差}：分子取\textbf{下界}：$\arctan x<\dfrac\pi2$，故
\[
2-\arctan x>2-\frac\pi2=\frac{4-\pi}{2}>0 .
\]
分母取\textbf{上界}：$\abs{1-x^2}=\abs{1+x}\cdot\abs{1-x}$。

\textbf{（3）关区间}：取 $\delta_1=\dfrac12$，则 $\dfrac32<x<\dfrac52$，故 $\abs{1+x}<\dfrac52$，于是
\[
\abs{1-x^2}<\frac52\abs{x-1}.
\]
\textbf{（4）放缩}：
\[
\abs{\frac{2-\arctan x}{1-x^2}}>\frac{(4-\pi)/2}{\frac52\abs{x-1}}=\frac{4-\pi}{5\abs{x-1}}.
\]
\textbf{（5）反解}：取 $\delta=\min\left\{\dfrac12,\dfrac{4-\pi}{5M}\right\}$，则 $0<\abs{x-1}<\delta$ 时上式 $>M$。$\blacksquare$
\end{solution}

\begin{notebox}[这道题的双侧极限其实不存在]
$\lim_{x\to1^-}f(x)=+\infty$，$\lim_{x\to1^+}f(x)=-\infty$（分子恒正、分母在右侧为负），所以 $\lim_{x\to1}f(x)$ 在“有限值”意义下不存在；题目与官方答案采用的是“$\abs{f}\to\infty$ 即称无穷大”的口径。\textbf{卷面上只能写“$\abs{f(x)}\to\infty$，即 $f$ 为无穷大”，不要写“$\lim f=+\infty$”。}
\end{notebox}

\begin{warnbox}[官方答案有错 · 2]
\textbf{2024--2025 卷第 11 题}：证明当 $x\to0$ 时 $\dfrac{x}{\sin^2x+x^3}$ 为无穷大。

官方答案先取 $\abs{x}<\frac12=\delta_1$，然后
\[
\abs{\frac{x}{\sin^2x+x^3}}\ge\frac{\abs{x}}{\abs{x}^2+\abs{x}^3}
=\frac1{\abs{x}(1+\abs{x})}\ \ge\ \frac1{\abs{x}}>G ,
\]
得 $\delta_2=\dfrac1G$。\textbf{最后一步方向反了}：正确的方向是
\[
\frac1{\abs{x}(1+\abs{x})}\ \le\ \frac1{\abs{x}},
\]
因为 $1+\abs{x}\ge1$。只有 $\abs{x}\le0$ 时 $\ge\frac1{\abs{x}}$ 才成立，而那是空话。

\textbf{反例}：取 $G=10$，官方 $\delta=0.1$；取 $x=0.09631$（满足 $0<\abs{x}<0.1$），实算 $\abs{f(x)}\approx9.497<10=G$，反例成立。该区间上 $\abs{f}$ 的最小值约 $9.12$（在 $x=0.1$ 处），同样不满足。
\end{warnbox}

\begin{exa}[2024--2025 卷第 11 题（订正版）]
证明当 $x\to0$ 时 $\dfrac{x}{\sin^2x+x^3}$ 为无穷大。
\end{exa}
\begin{solution}
\textbf{（1）写定义}：任给 $M>0$，需找 $\delta>0$，使 $0<\abs{x}<\delta$ 时 $\abs{\dfrac{x}{\sin^2x+x^3}}>M$。

\textbf{（2）化差}：用 $\abs{\sin x}\le\abs{x}$ 与三角不等式，
\[
\abs{\frac{x}{\sin^2x+x^3}}\ge\frac{\abs{x}}{\abs{\sin x}^2+\abs{x}^3}
\ge\frac{\abs{x}}{\abs{x}^2+\abs{x}^3}.
\]
\textbf{（3）关区间}：取 $\delta_1=\dfrac12$，则 $1+\abs{x}\le\dfrac32<2$，于是
\[
\frac{\abs{x}}{\abs{x}^2+\abs{x}^3}=\frac1{\abs{x}(1+\abs{x})}>\frac1{2\abs{x}}.
\]
（注意这里把 $1+\abs{x}$ 换成上界 $2$，正好符合“求下界时分母取上界”。）

\textbf{（4）放缩}：$\abs{f(x)}>\dfrac1{2\abs{x}}$。\\[2pt]
\textbf{（5）反解}：取 $\delta=\min\left\{\dfrac12,\dfrac1{2M}\right\}$，则 $0<\abs{x}<\delta$ 时 $\abs{f(x)}>\dfrac1{2\abs{x}}>\dfrac1{2\delta}\ge M$。$\blacksquare$
\end{solution}

\begin{methodbox}[这一类“求下界”的标准套路]
$\dfrac{\abs{x}}{\abs{x}^2+\abs{x}^3}$ 本身已经比 $\dfrac1{\abs{x}}$ \textbf{小}，所以\textbf{不能}用 $\ge\dfrac1{\abs{x}}$；正确的下界是 $\dfrac1{2\abs{x}}$（因为 $1+\abs{x}\le2$）。写放缩链时每一步都问一句：“这一步真的是 $\ge$ 吗？”
\end{methodbox}

% ============================================================
\section{十条高频错误对照}
% ============================================================

\begin{enumerate}
\item \no\ 直接写“取 $\delta=\eps/2$”，不写“任给 $\eps>0$”。
      \yes\ 先写定义句式，再进入求解。

\item \no\ $\delta$ 里出现 $x$（如 $\delta=\eps\abs{x-1}$）。\quad
      \yes\ $\delta$ 只能依赖 $\eps$ 和题中常数。

\item \no\ 把 $\abs{f-A}$ 往小了估（$\le$ 写成 $\ge$）。\quad
      \yes\ 求“$\abs{f-A}<\eps$”要\textbf{估上界}；只有无穷大型才估下界。

\item \no\ 忘了 $0<\abs{x-x_0}$（把 $x=x_0$ 也算进去）。\quad
      \yes\ 根式差、分式差在 $x=x_0$ 处可能无定义，必须去心。

\item \no\ 省掉 $\delta_1$，直接用全局常数。\quad
      \yes\ 先取 $\delta_1$ 定住分母/根号，再反解 $\eps/C$。

\item \no\ 最后写 $\delta=\eps/C$，漏掉 $\min$。\quad
      \yes\ $\delta=\min\{\delta_1,\eps/C\}$；$X$ 型对应 $\max$。

\item \no\ $x\to\infty$ 时仍写 $\delta$。\quad
      \yes\ 模板 B 的阈值记作 $X$，条件写 $\abs{x}>X$ 或 $x>X$。

\item \no\ 无穷大型把 $\abs{f}>M$ 写成 $f<M$，或只证单侧。\quad
      \yes\ 若结论是 $f\to+\infty$，必须双侧同号趋正；否则只能写 $\abs{f}\to\infty$。

\item \no\ 含参题试图“解出 $a$ 的范围”。\quad
      \yes\ 用 $\abs{a}$ 把参数压成常数，再取 $X_1$ 让高阶项占优。

\item \no\ 用极限四则运算“证明”定义（如“因为 $\lim x=1$，$\lim\sqrt{x+3}=2$，所以乘积为 $2$”）。\quad
      \yes\ 定义证明不能用待证的极限性质，必须回到 $\eps$--$\delta$ 的放缩。
\end{enumerate}

% ============================================================
\section{练习（附答案）}
% ============================================================

\begin{enumerate}
\item 用定义证明 $\displaystyle\lim_{x\to3}(2x+1)=7$。
\item 用定义证明 $\displaystyle\lim_{x\to1}\frac{x^2-1}{x-1}=2$。
\item 用定义证明 $\displaystyle\lim_{x\to0}x\sin\frac1x=0$。
\item 用定义证明 $\displaystyle\lim_{x\to+\infty}\frac{3x+1}{x-2}=3$。
\item 用无穷大定义证明 $\displaystyle\lim_{x\to2}\frac{1}{(x-2)^2}=+\infty$。
\end{enumerate}

\begin{keypoint}[答案与关键步骤]
\begin{enumerate}
\item $\abs{2x+1-7}=2\abs{x-3}$，取 $\delta=\dfrac{\eps}{2}$（无需 $\delta_1$）。
\item 去心邻域内 $\dfrac{x^2-1}{x-1}=x+1$，故 $\abs{f-2}=\abs{x-1}$，取 $\delta=\eps$。
\item $\abs{x\sin\frac1x}\le\abs{x}$，取 $\delta=\eps$。
\item $\abs{\dfrac{3x+1}{x-2}-3}=\dfrac{7}{x-2}$；$x\ge4$ 时 $x-2\ge\dfrac x2$，故 $\le\dfrac{14}{x}$，取 $X=\max\left\{4,\dfrac{14}{\eps}\right\}$。
\item 因 $(x-2)^2>0$（双侧同号），$\dfrac1{(x-2)^2}>M\iff\abs{x-2}<\dfrac1{\sqrt M}$，取 $\delta=\dfrac1{\sqrt M}$。这是\textbf{真正的} $+\infty$（不需要写 $\abs{f}$）。
\end{enumerate}
\end{keypoint}

\begin{warnbox}[再加一道：否定式]
证明 $\lim_{x\to0}\cos\dfrac1x$ 不存在。\\
提示：取 $\eps_0=\dfrac12$，用 $x_n=\dfrac1{2n\pi}$（$\cos=1$）与 $y_n=\dfrac1{(2n+1)\pi}$（$\cos=-1$），论证与第 2.6 节的 $\sin\frac1x$ 完全同构。
\end{warnbox}

% ============================================================
\section{考场执行清单（写完扫 20 秒）}
% ============================================================

\begin{enumerate}
\item[$\square$] 结论里 $x\to$ 的是 $x_0$ 还是 $\infty$？$\infty$ 是“有限值极限”还是“无穷大”？
\item[$\square$] 定义句式写全了吗（$\forall$ / $\exists$ / 范围 / 不等式）？
\item[$\square$] 有没有先用 $\delta_1$ / $X_1$ 把 $x$ 关进小区间？
\item[$\square$] 分母拿到的是上界（求下界时）还是下界（求上界时）？方向对吗？
\item[$\square$] $\delta$ / $X$ 是 $\min$ / $\max$ 形式吗（两个条件同时满足）？
\item[$\square$] 含参题：参数被常数控住了吗？
\item[$\square$] 代一个具体值抽查过吗（$M=10$ 或 $\eps=0.1$）？
\end{enumerate}

% ============================================================
\appendix
\section{17 套真题第 11 题全表（2008--2025）}
% ============================================================

\begin{notebox}[数据来源与核对方式]
15 套卷面 PDF 用 pymupdf 定位第 11 题所在行后 4 倍渲染逐字读取；2012--2013、2014--2015 两套\textbf{只有答案 PDF}（卷面缺失），从答案页定位读取；2024--2025 从答案页读取。
表中“关键取法”一列是本讲义按五步骨架重新推出的常数，已用数值抽样复核（脚本见 \texttt{Homework/\_review/tools/verify\_handout\_examples.py}）。
\end{notebox}

{\footnotesize
\renewcommand{\arraystretch}{1.25}
\begin{longtable}{@{}>{\raggedright\arraybackslash}p{1.65cm}>{\raggedright\arraybackslash}p{2.5cm}>{\raggedright\arraybackslash}p{7.3cm}>{\raggedright\arraybackslash}p{3.85cm}@{}}
\toprule
学年 & 形式 & 题面 & 关键取法 \\
\midrule
\endfirsthead
\toprule
学年 & 形式 & 题面 & 关键取法 \\
\midrule
\endhead
\bottomrule
\endfoot
2008--2009 & A $\eps$--$\delta$ & $\lim_{x\to0}\dfrac{x-1}{x^2+1}=-1$ & $\delta=\min\{1,\eps/2\}$ \\
2009--2010 & C 无穷大（$x\to+\infty$） & $\lim_{x\to+\infty}(x^3-4x)=+\infty$ & $X=\max\{3,(2M)^{1/3}\}$ \\
2010--2011 & B $\eps$--$X$ & $\lim_{x\to+\infty}\dfrac{x^2+x}{x^2-2}=1$ & $X=\max\{2,4/\eps\}$ \\
2011--2012 & A $\eps$--$\delta$ & $\lim_{x\to2}\dfrac{1}{x-1}=1$ & $\delta=\min\{1/2,\eps/2\}$ \\
2012--2013$^\dagger$ & A $\eps$--$\delta$ & $\lim_{x\to2}\dfrac{1}{1+x^2}=\dfrac15$ & $\delta=\min\{1,\eps\}$ \\
2013--2014 & B $\eps$--$X$ & $\lim_{x\to+\infty}\dfrac{\sqrt{2x+\sqrt x}-\sqrt x}{1+x}=0$ & $X=\max\{1,4/\eps^2\}$ \\
2014--2015$^\dagger$ & A $\eps$--$\delta$ & $\lim_{x\to0}\left(\dfrac{1}{1+x}-\sin x\right)=1$ & $\delta=\min\{1/2,\eps/3\}$ \\
2015--2016 & C 无穷大（$x\to1$） & $\lim_{x\to1}\dfrac{2-\arctan x}{1-x^2}=\infty$ & $\delta=\min\{1/2,(4-\pi)/(5M)\}$ \\
2016--2017 & B $\eps$--$X$ & $\lim_{x\to\infty}\dfrac{2x^2}{x^2-3x+4}=2$ & $X=\max\{8,14/\eps\}$ \\
2017--2018 & A $\eps$--$\delta$ & $\lim_{x\to1}x\sqrt{x+3}=2$ & $\delta=\min\{1,\eps/3\}$ \\
2018--2019 & B $\eps$--$X$（含参 $a$） & $\lim_{x\to+\infty}\dfrac{x^2+\ee^{-x}}{x^3-ax+1}=0$ & $X=\max\{1,\sqrt{2(\abs{a}+1)},4/\eps\}$ \\
2019--2020 & B $\eps$--$X$ & $\lim_{x\to\infty}\dfrac{x+1}{x^2+2}=0$ & $X=\max\{1,2/\eps\}$ \\
2020--2021 & A $\eps$--$\delta$ & $\lim_{x\to\pi/2}x\sin x=\dfrac\pi2$ & $\delta=\min\{1,4\eps/(4+\pi)\}$ \\
2021--2022 & A $\eps$--$\delta$ & $\lim_{x\to1}\dfrac{\sqrt{3-x}-\sqrt{x+1}}{x-1}=-\dfrac{1}{\sqrt2}$ & $\delta=\min\{1/2,5\eps/2\}$ \\
2022--2023 & A $\eps$--$\delta$ & $\lim_{x\to1}\dfrac{\sqrt x}{2x-1}=1$ & $\delta=\min\{1/4,\eps/4\}$ \\
2023--2024 & B $\eps$--$X$ & $\lim_{x\to\infty}\dfrac{2x^2+x+\arctan x}{x^2-1}=2$ & $X=\max\{4,4/\eps\}$ \\
2024--2025 & C 无穷大（$x\to0$） & $x\to0$ 时 $\dfrac{x}{\sin^2x+x^3}$ 为无穷大 & $\delta=\min\{1/2,1/(2M)\}$ \\
\end{longtable}}

\noindent{\footnotesize $^\dagger$ 该学年只有参考答案 PDF，卷面缺失；题面据答案页转述。}

\begin{center}
\small
\begin{tabular}{@{}lccc@{}}
\toprule
形式 & A $\eps$--$\delta$ & B $\eps$--$X$ & C 无穷大 $M$ 型 \\
\midrule
套数 & 8 & 6 & 3 \\
建议练习配比 & 3 & 2 & 1 \\
\bottomrule
\end{tabular}
\end{center}

\noindent 三种形式\textbf{接近均匀轮换}，不要只练顺手的一种；模板 C 只考 3 次，却年年有人栽（放缩方向反了）。

% ============================================================
\section{考前 20 分钟速查}
% ============================================================

\begin{keypoint}[三行定义]
\begin{align*}
\text{A：}& \forall\eps>0\ \exists\delta>0:\ 0<\abs{x-x_0}<\delta\Rightarrow\abs{f-A}<\eps;\\
\text{B：}& \forall\eps>0\ \exists X>0:\ x>X\Rightarrow\abs{f-A}<\eps;\\
\text{C：}& \forall M>0\ \exists\delta>0:\ 0<\abs{x-x_0}<\delta\Rightarrow\abs{f}>M .
\end{align*}
\end{keypoint}

\begin{methodbox}[五步]
写定义 $\to$ 化差（有理化/通分/拆项）$\to$ 先用 $\delta_1$ 关区间 $\to$ 放缩成 $C\abs{x-x_0}$ 或 $C/\abs{x}$ 或 $k/\abs{x-x_0}$ $\to$ 反解（$\min$ / $\max$）。
\end{methodbox}

\begin{warnbox}[三条最容易丢分的]
\begin{enumerate}
\item 方向：求上界（$\abs{f-A}<\eps$）就\textbf{分母取下界、分子取上界}；求下界（$\abs{f}>M$）就反过来。
\item 忘 $\delta_1$ 或忘 $\min$。\quad
\item $f\to+\infty$ 要求双侧同号；否则只能写 $\abs{f}\to\infty$。
\end{enumerate}
\end{warnbox}

\end{document}
